
\documentclass[11pt]{article}

\usepackage[margin=0.9in]{geometry}
\usepackage{amsmath,amssymb}
\usepackage{array}
\usepackage{booktabs}
\usepackage{enumitem}
\usepackage{listings}
\usepackage{xcolor}
\usepackage{tikz}
\usepackage{hyperref}

\definecolor{backcolour}{rgb}{0.95,0.95,0.92}
\definecolor{codeblue}{rgb}{0.05,0.15,0.55}
\definecolor{codegreen}{rgb}{0.0,0.45,0.0}

\lstset{
    backgroundcolor=\color{backcolour},
    basicstyle=\ttfamily\small,
    keywordstyle=\color{codeblue}\bfseries,
    commentstyle=\color{codegreen},
    stringstyle=\color{purple},
    breaklines=true,
    frame=single,
    showstringspaces=false,
    tabsize=4
}

\title{CPSC 250 \\ Big Program Two: Credit Card Debt Search Program}
\author{Edward Brash}
\date{}

\begin{document}

\maketitle

\section{Introduction}

In this project, we consider a realistic programming task involving a large CSV data file.

Suppose we are given a large CSV file containing information about individuals in the United States, including:

\begin{itemize}
    \item last names,
    \item states,
    \item credit card debt values.
\end{itemize}

Our goal is to write a Python program that can search this data efficiently and answer questions such as:

\begin{enumerate}
    \item Find all people whose names start with a particular string.
    \item Find all people with debt below some limit.
    \item Find all people from a particular state.
    \item Find the person with the highest debt.
\end{enumerate}

Although this may initially seem like a large and intimidating task, the key idea is:

\begin{quote}
Break the large problem into many smaller problems.
\end{quote}

Each smaller problem is manageable and usually only requires a few lines of code.

\section{Questions We Must Answer}

Before writing code, we should think carefully about the design of the program.

Some important questions are:

\begin{enumerate}
    \item What does the CSV file look like?
    \item What data structures should we use?
    \item What algorithms are needed for searching?
\end{enumerate}

These are examples of software design questions.

\section{Understanding the CSV File}

Suppose our CSV file is called:

\begin{lstlisting}
customer_data.csv
\end{lstlisting}

and contains rows such as:

\begin{lstlisting}
Smith,VA,2450.75
Johnson,CA,10234.10
Jones,NY,500.00
Brown,TX,7890.50
\end{lstlisting}

Each row contains three pieces of information:

\begin{enumerate}
    \item Last name,
    \item State,
    \item Credit card debt.
\end{enumerate}

Notice that:

\begin{itemize}
    \item names are strings,
    \item states are strings,
    \item debts are floating point numbers.
\end{itemize}

\section{Choosing a Data Structure}

The notes propose a very simple first design idea:

\begin{quote}
Use three ordered lists.
\end{quote}

Specifically:

\begin{lstlisting}[language=Python]
names = []
states = []
debts = []
\end{lstlisting}

These are called \textbf{parallel lists}.

\subsection{Why Parallel Lists Work}

The important idea is that the entries at the same index correspond to the same person.

For example:

\begin{center}
\begin{tabular}{cccc}
\toprule
Index & names & states & debts \\
\midrule
0 & Smith & VA & 2450.75 \\
1 & Johnson & CA & 10234.10 \\
2 & Jones & NY & 500.00 \\
3 & Brown & TX & 7890.50 \\
\bottomrule
\end{tabular}
\end{center}

Thus:

\begin{lstlisting}[language=Python]
names[2]
\end{lstlisting}

corresponds to:

\begin{lstlisting}[language=Python]
states[2]
\end{lstlisting}

and:

\begin{lstlisting}[language=Python]
debts[2]
\end{lstlisting}

because all three refer to the same customer.

\section{Step 1: Read the Data from the File}

The first major task is to read the CSV file and store the information in the lists.

The notes propose writing a function:

\begin{lstlisting}[language=Python]
def read_customer_data(filename):
\end{lstlisting}

This function should:

\begin{enumerate}
    \item open the file,
    \item read every row,
    \item extract the data,
    \item store the data in the lists,
    \item return the lists.
\end{enumerate}

\section{Opening a CSV File}

Python provides built-in file handling tools.

A common pattern is:

\begin{lstlisting}[language=Python]
with open(filename, "r") as file:
\end{lstlisting}

Important ideas:

\begin{itemize}
    \item \texttt{"r"} means ``read mode'',
    \item \texttt{file} becomes a file object,
    \item the \texttt{with} statement automatically closes the file.
\end{itemize}

\section{Using the CSV Module}

Python includes a built-in module called \texttt{csv} for working with CSV files.

We begin by importing it:

\begin{lstlisting}[language=Python]
import csv
\end{lstlisting}

Then we create a CSV reader:

\begin{lstlisting}[language=Python]
reader = csv.reader(file)
\end{lstlisting}

The reader behaves like a loop over rows.

Each row is itself a list of strings.

For example, a row might look like:

\begin{lstlisting}[language=Python]
["Smith", "VA", "2450.75"]
\end{lstlisting}

\section{Reading Rows from the File}

We loop through the rows:

\begin{lstlisting}[language=Python]
for row in reader:
\end{lstlisting}

Inside the loop:

\begin{itemize}
    \item \texttt{row[0]} is the name,
    \item \texttt{row[1]} is the state,
    \item \texttt{row[2]} is the debt.
\end{itemize}

We append these values to the appropriate lists.

\section{The \texttt{append()} Method}

Python lists provide the method:

\begin{lstlisting}[language=Python]
my_list.append(value)
\end{lstlisting}

This adds a new value to the end of the list.

For example:

\begin{lstlisting}[language=Python]
names.append("Smith")
\end{lstlisting}

adds the string \texttt{"Smith"} to the end of the list.

\section{Converting Debt to a Float}

When data is read from a CSV file, everything initially appears as a string.

For example:

\begin{lstlisting}[language=Python]
row[2] == "2450.75"
\end{lstlisting}

But debt values should behave numerically.

Therefore we convert them:

\begin{lstlisting}[language=Python]
float(row[2])
\end{lstlisting}

This converts the string into a floating point number.

\section{Complete File Reading Function}

Putting everything together:

\begin{lstlisting}[language=Python]
import csv

def read_customer_data(filename):

    names = []
    states = []
    debts = []

    with open(filename, "r") as file:

        reader = csv.reader(file)

        for row in reader:

            names.append(row[0])
            states.append(row[1])
            debts.append(float(row[2]))

    return names, states, debts
\end{lstlisting}

\section{Returning Multiple Values}

Notice the line:

\begin{lstlisting}[language=Python]
return names, states, debts
\end{lstlisting}

Python allows functions to return multiple values.

We can capture them using:

\begin{lstlisting}[language=Python]
names, states, debts = read_customer_data("customer_data.csv")
\end{lstlisting}

This is extremely useful when working with related data structures.

\section{Step 2: Get Information from the User}

Next, the program asks the user for search information.

The notes suggest three kinds of input:

\begin{enumerate}
    \item a debt limit,
    \item a search phrase,
    \item a state abbreviation.
\end{enumerate}

For example:

\begin{lstlisting}[language=Python]
debt_limit = float(input("Enter debt limit: "))
search_phrase = input("Enter search phrase: ")
st = input("Enter state abbreviation: ")
\end{lstlisting}

Notice:

\begin{itemize}
    \item debt limits should be converted to numbers,
    \item names and states remain strings.
\end{itemize}

\section{Searching for Names}

Suppose we want all names beginning with:

\begin{lstlisting}
"Jo"
\end{lstlisting}

We loop through the list of names.

\begin{lstlisting}[language=Python]
for n in range(len(names)):
\end{lstlisting}

The expression:

\begin{lstlisting}[language=Python]
range(len(names))
\end{lstlisting}

produces indices:

\[
0, 1, 2, 3, \ldots
\]

Using indices is important because we need access to the matching state and debt information.

\section{The \texttt{startswith()} String Method}

Python strings provide a method called:

\begin{lstlisting}[language=Python]
startswith()
\end{lstlisting}

Example:

\begin{lstlisting}[language=Python]
if names[n].startswith(search_phrase):
\end{lstlisting}

This checks whether the string begins with the desired text.

For example:

\begin{lstlisting}[language=Python]
"Johnson".startswith("Jo")
\end{lstlisting}

returns:

\begin{lstlisting}
True
\end{lstlisting}

while:

\begin{lstlisting}[language=Python]
"Smith".startswith("Jo")
\end{lstlisting}

returns:

\begin{lstlisting}
False
\end{lstlisting}

\section{Complete Name Search Example}

\begin{lstlisting}[language=Python]
print("Matching names:")

for n in range(len(names)):

    if names[n].startswith(search_phrase):

        print(names[n], states[n], debts[n])
\end{lstlisting}

\section{Searching by Debt}

Suppose we want everyone with debt less than \$2000.

We compare each debt value against the limit.

\begin{lstlisting}[language=Python]
print("People below debt limit:")

for n in range(len(debts)):

    if debts[n] < debt_limit:

        print(names[n], states[n], debts[n])
\end{lstlisting}

Notice that this is a simple linear search through the data.

\section{Searching by State}

Suppose we want all customers from Virginia.

\begin{lstlisting}[language=Python]
print("People from state:", st)

for n in range(len(states)):

    if states[n] == st:

        print(names[n], states[n], debts[n])
\end{lstlisting}

This uses the equality operator:

\begin{lstlisting}[language=Python]
==
\end{lstlisting}

to compare strings.

\section{Step 3: Find the Highest Credit Card Debt}

Another important task is finding the maximum debt value.

The notes use a classic algorithm based on tracking the index of the largest value found so far.

\section{Understanding the Maximum Search Algorithm}

The algorithm is:

\begin{enumerate}
    \item Assume the first value is the maximum.
    \item Loop through the remaining values.
    \item If a larger value is found, update the index.
    \item At the end, the stored index corresponds to the largest debt.
\end{enumerate}

\section{Implementing the Maximum Search}

\begin{lstlisting}[language=Python]
index_max = 0

for n in range(len(debts)):

    if debts[n] > debts[index_max]:

        index_max = n
\end{lstlisting}

Important idea:

\begin{quote}
We store the index of the maximum value, not the value itself.
\end{quote}

This is useful because the index also gives access to the matching name and state.

\section{Printing the Maximum Debt Customer}

After the loop finishes:

\begin{lstlisting}[language=Python]
print("Highest debt customer:")
print(names[index_max], states[index_max], debts[index_max])
\end{lstlisting}

Suppose:

\begin{center}
\begin{tabular}{cccc}
\toprule
Index & Name & State & Debt \\
\midrule
0 & Smith & VA & 2450 \\
1 & Johnson & CA & 10234 \\
2 & Jones & NY & 500 \\
\bottomrule
\end{tabular}
\end{center}

Then:

\begin{lstlisting}[language=Python]
index_max == 1
\end{lstlisting}

because Johnson has the largest debt.

\section{Building Programs Incrementally}

A very important lesson from this project is:

\begin{quote}
Large programs should be built incrementally.
\end{quote}

Instead of trying to write everything at once:

\begin{enumerate}
    \item Write one small piece.
    \item Test it.
    \item Move to the next piece.
    \item Test again.
\end{enumerate}

This is how real software development works.

\section{Program Decomposition}

At first, a 100-line program can appear overwhelming.

However:

\begin{quote}
Most large programs are really collections of many small programs.
\end{quote}

In this project:

\begin{itemize}
    \item reading the file is only a few lines,
    \item searching names is only a few lines,
    \item searching debts is only a few lines,
    \item finding a maximum is only a few lines.
\end{itemize}

The complexity comes from combining many simple ideas together.

\section{Complete Example Program}

\begin{lstlisting}[language=Python]
import csv


def read_customer_data(filename):

    names = []
    states = []
    debts = []

    with open(filename, "r") as file:

        reader = csv.reader(file)

        for row in reader:

            names.append(row[0])
            states.append(row[1])
            debts.append(float(row[2]))

    return names, states, debts


# Read data from the CSV file
names, states, debts = read_customer_data("customer_data.csv")


# Get information from the user
debt_limit = float(input("Enter debt limit: "))
search_phrase = input("Enter search phrase: ")
st = input("Enter state abbreviation: ")


# Search for names
print()
print("Names beginning with", search_phrase)

for n in range(len(names)):

    if names[n].startswith(search_phrase):

        print(names[n], states[n], debts[n])


# Search by debt
print()
print("People below debt limit")

for n in range(len(debts)):

    if debts[n] < debt_limit:

        print(names[n], states[n], debts[n])


# Search by state
print()
print("People from", st)

for n in range(len(states)):

    if states[n] == st:

        print(names[n], states[n], debts[n])


# Find highest debt
index_max = 0

for n in range(len(debts)):

    if debts[n] > debts[index_max]:

        index_max = n


print()
print("Highest debt customer:")

print(names[index_max], states[index_max], debts[index_max])
\end{lstlisting}

\section{Possible Improvements}

This design works well for learning purposes, but there are many possible improvements.

For example:

\begin{itemize}
    \item use dictionaries instead of parallel lists,
    \item use classes and objects,
    \item use pandas,
    \item sort the data,
    \item use binary search,
    \item add error checking,
    \item handle malformed CSV files,
    \item make searches case-insensitive.
\end{itemize}

However, the simple version is ideal for learning the fundamentals.

\section{Computational Complexity}

All searches in this program are linear searches.

That means:

\[
\mathcal{O}(n)
\]

time complexity.

Why?

Because in the worst case, we may need to examine every entry in the list.

For a file with:

\[
10{,}000
\]

rows, this is usually perfectly acceptable.

For millions of rows, more sophisticated data structures might become necessary.

\section{Key Takeaways}

\begin{itemize}
    \item CSV files store structured tabular data.
    \item The \texttt{csv} module helps read CSV files safely.
    \item Parallel lists are a simple but useful data structure.
    \item File data is initially read as strings.
    \item Numerical values often require conversion using \texttt{float()} or \texttt{int()}.
    \item Linear searches are straightforward to implement.
    \item The \texttt{startswith()} method is useful for string searching.
    \item Maximum search algorithms often track indices rather than values.
    \item Large programs become manageable when broken into smaller pieces.
\end{itemize}

\end{document}
